\section{The Server}\label{sec:design}

\oneport{} is the paper's own minimal server, written in C++23 and built with CMake. This section
describes it as the frozen design specifies it and as the code implements it. It holds exactly the
features the hypotheses study, and nothing more.

\subsection{Structure}\label{sec:structure}

One binary takes four flags: the mode (one-port, dedicated or stub), the detection mode (replay or
peek), the dispatch mode (in-process or relay) and the backend (epoll, io\_uring or IOCP). A
process runs one backend. It runs one worker thread per server core; the primary cells use one
core. The event loop of each worker derives from the minimal loop of an earlier study by the same
author~\cite{tsvetanov2026wake}. From it come the io\_uring ring on raw system calls, created for
a single issuer with deferred task work and enabled on the worker. So do the rules for memory that
the kernel fills, under MemorySanitizer, the stop paths and the IOCP port setup. Sockets, timers,
detection, handlers and counters are new.

A connection's state is a tagged union of handler states, and dispatch is a switch on the tag,
with no virtual call. A dedicated listener sets the tag at accept; a one-port listener sets it
when detection classifies the connection. From then on both modes run the same code. The mode
flag is read once, by the listener setup. A structural test fails if any other source file reads
it, so the handlers cannot differ between the modes.

\subsection{Listeners}\label{sec:listeners}

In one-port mode a process has one listening TCP socket. When several workers serve it, they share
it: on epoll each worker's set holds it with \texttt{EPOLLEXCLUSIVE}, and on io\_uring each
worker's ring holds a multishot accept on it. On IOCP every listener keeps \NAcceptex{}
\texttt{AcceptEx} requests outstanding, without a receive buffer, in both
modes~\cite{microsoft2026acceptex}. A group of \texttt{SO\_REUSEPORT} sockets, one per worker, is
a second layout that the binary also builds. On IOCP the server runs one worker and refuses more,
so the secondary cell with two server cores on IOCP was not run (Section~\ref{sec:res-secondary}).
In dedicated mode a process opens one listening socket per
protocol (HTTP/1.1, h2c, TLS, MQTT, SSH and SMTP) on consecutive ports. It has no detection, no
detection timer and no detection budget. The backend, the workers, the handlers, the buffers and
the PROXY setting are the same in both modes. That is the whole difference between them, and it
is what the cost family measures.

\subsection{The Detection Table}\label{sec:table}

Each matcher is a \texttt{constexpr} descriptor. Its function maps the bytes received so far to
one of three answers: yes, no, or more needed. It also states the fewest and the most bytes it
needs to say yes, and the set of first bytes it can accept. The set is derived from the function, so the two cannot
disagree. Table~\ref{tab:matchers} gives the matchers as frozen. The methods of HTTP/1.1 are the
eight of HTTP semantics~\cite{fielding2022rfc9110} and PATCH~\cite{dusseault2010rfc5789}.

\begin{table}[H]
\caption{The detection table, as frozen. The PROXY parser is not part of the table; it runs before
it, on a listener configured for PROXY only.\label{tab:matchers}}
\centering\small
\begin{tabularx}{\textwidth}{@{}l X l@{}}
\toprule
\textbf{Matcher} & \textbf{Accepts} & \textbf{Decides at} \\
\midrule
TLS & first byte \texttt{0x16}; second byte \texttt{0x03}; third byte any; a record length of at
most \TlsRecordMax{}; sixth byte \texttt{0x01}, a ClientHello~\cite{rescorla2018rfc8446} &
\TlsNeed{} bytes \\
h2c & the preface of \HtwocNeed{} bytes, exactly~\cite{thomson2022rfc9113} & \HtwocNeed{} bytes,
or the first byte that differs \\
HTTP/1.1 & at most one leading empty line, then one of the nine methods, then exactly one
space~\cite{fielding2022rfc9112} & at most \HttpNeedMax{} bytes \\
SSH & \texttt{SSH-}~\cite{ylonen2006rfc4253} & \SshNeed{} bytes \\
MQTT & first byte \texttt{0x10}, a CONNECT; a Remaining Length of \MqttRlMin{} to \MqttRlMax{}
bytes; then the protocol name \texttt{MQTT} with its length, and the level of MQTT~3.1.1 or of
MQTT~5.0~\cite{banks2014mqtt311,banks2019mqtt5} & \MqttNeedMin{} to \MqttNeedMax{} bytes \\
PROXY v1, v2 & v2 needs \ProxyVtwoNeed{} bytes with the first \ProxyVtwoMatch{} matching; v1
needs \ProxyVoneNeed{} with the first \ProxyVoneMatch{} equal to \texttt{PROXY}; a prefix that
cannot match closes the connection; a v2 header longer than \ProxyVtwoMax{} bytes is
rejected~\cite{tarreau2026proxy} & then the whole header \\
\bottomrule
\end{tabularx}
\end{table}

From the table the compiler builds an array of all byte values that maps a first byte to its
candidate matchers. Within a bucket, the matcher with the smaller need runs first, for speed only.
Checks at compile time run every matcher over a corpus of samples built from the
specifications' grammars. Each sample is accepted by its own matcher. No matcher accepts a prefix
of another protocol's sample. Running each bucket in reverse order changes no decision on the
corpus. The first-byte sets equal those of the frozen table. A build that broke any of these
would not compile, so on the corpus the classes are prefix-disjoint and the order of the table
cannot change a decision. The test suite runs the same corpus again at run time, split at every
position.

Classification keeps no state between calls. The bytes stay where the detection mode keeps them
(Section~\ref{sec:modes}), and the table runs on all of them each time more arrive. So a signature
fed in any split gets the decision it gets in one write. An empty candidate set after the first
byte rejects the connection at once. Two budgets are compile-time constants. An undecided
connection may hold at most \Bdec{} = \BDec{} bytes after any PROXY header. A ClientHello
reassembled in pass-through may hold at most \Bch{} = \BCh{} bytes. A connection that exceeds a budget while it
is undecided is rejected. On a listener configured for PROXY, the header is parsed first and
consumed exactly, so the table sees the first application byte. That order removes the overlap of
the binary PROXY header with HTTP/1.1's leading empty line.

\subsection{Detection Modes}\label{sec:modes}

\runin{Replay.} The first read goes into the buffer the handler would read into in dedicated mode.
The matchers run on it, and the handler then parses the same buffer from its start. No byte is
copied for the replay.

\runin{Peek.} The bytes are copied with \texttt{MSG\_PEEK} into a scratch buffer of the worker and
stay in the kernel. The handler later reads them as in dedicated mode. An undecided peek sets
\texttt{SO\_RCVLOWAT} to the fewest bytes in all with which some candidate could still say yes.
The socket then reports readiness again only when that many bytes are queued~\cite{man2026socket}.
Before the handler reads, one more call sets the mark back to one byte.

The backends realise the two modes differently.
\begin{itemize}
\item \textbf{epoll.} Every connection is registered edge-triggered with \texttt{EPOLLRDHUP}, in
both modes~\cite{man2026epoll}. A readiness event is followed by a synchronous receive (replay) or
peek. \texttt{EPOLLRDHUP} reports a peer that shut down writing, which a peek alone cannot show
while bytes are queued.
\item \textbf{io\_uring.} In replay, a read on the socket is posted with a buffer that the kernel
selects from a ring of provided buffers when data arrives~\cite{man2026iouring}. A pending read
therefore holds no buffer of its own. The opcode is \texttt{IORING\_OP\_READ}, a reading the
revision log records. In peek, the worker posts a poll for readable data or a half-close, then
peeks synchronously. A test pins that a poll armed again after an undecided peek completes only
at the low-water mark.
\item \textbf{IOCP.} A zero-byte \texttt{WSARecv} serves as the readiness signal, followed by a
synchronous receive or peek on the non-blocking socket~\cite{microsoft2026wsarecv}. Windows has
no \texttt{SO\_RCVLOWAT}, so nothing can wait in the kernel for more bytes. An undecided peek
therefore switches the connection to replay: the queued bytes are read into the handler's buffer
at once and matching continues there. The counters record each switch.
\end{itemize}
Each backend has one default detection mode, used by every cost and robustness cell. The proposed
default is replay on all three backends. Rule~E (Section~\ref{sec:order}) may change it per
backend after the pilot.

\runin{Memory of a pending connection.} By design, a pending connection holds its state, which is
fixed in size, and no data buffer while no byte has arrived, on every backend. With replay it
holds the handler's buffer from the first byte on; with peek, nothing in user space. In
pass-through it holds the partial ClientHello, at most \Bch{}. The test suite audits these bounds. IOCP has a second
receive form that posts the handler's buffer with the first \texttt{WSARecv}. It holds a buffer
while no byte has arrived, so it breaks the bound; it is built only for one secondary cell,
reported with that violation.

\subsection{Timers and the Decision Deadline}\label{sec:timers}

Three timers bound detection. The silence timer \Tfb, the decision deadline \Tdec{} and the PROXY
header deadline \Thdr{} are each \TimerS{}~s, a design choice. The PROXY specification asks for at
least that long~\cite{tarreau2026proxy}, and it is caddy-l4's default matching
timeout~\cite{caddyl42026}. \Tfb{} and \Tdec{} start together, at accept, or on a PROXY listener
when the header is complete. A listener may name one fallback protocol. The frozen rules are these.
\begin{enumerate}[label=(\alph*)]
\item A connection that has delivered no application byte when \Tfb{} expires goes to the fallback
handler. It goes in the first loop pass that handles the expiry, never earlier, unless the check
of (b) finds a byte.
\item Every loop pass handles all of its completions and readiness events before any expiry.
Before every fallback dispatch the server checks for a byte without waiting. Where no receive
was posted with a buffer it peeks one byte with \texttt{MSG\_PEEK}; where one was, it reaps
completions without waiting. If it finds a byte, the bytes alone decide the connection.
\item After a fallback dispatch, the byte transcript in both directions equals that of the same
client against the fallback's dedicated port. Only the timing differs.
\item A connection with at least one byte that no matcher has decided at \Tdec{} is closed as
undecided. It never goes to the fallback.
\item A connection whose bytes rule out every matcher is closed at once as rejected.
\item A connection that is silent at \Tdec{}, on a listener without a fallback, is closed as silent.
\item A client whose first byte is delivered after \Tfb{} goes to the fallback. That is the defined
price of the boundary, not an error.
\item A connection whose peer shuts down writing before a decision is closed at once: undecided
if it delivered a byte, silent if not.
\end{enumerate}
In relay dispatch a TLS connection is classified at its sixth byte, but its route needs the whole
ClientHello. \Tdec{} stays armed until the route is chosen, so it bounds that wait too.

\runin{The deadline queue.} Every pending connection of a listener waits under the same
durations, so its deadlines arrive in accept order. One intrusive first-in, first-out list per
listener and timer therefore keeps them sorted, with constant-time arming, cancelling and expiry.
The loop bounds each wait by the earliest head deadline: \EpollWaitTwo{} with a time
specification, the io\_uring wait with its timeout, and \texttt{GetQueuedCompletionStatusEx} in
whole milliseconds. After every wait the loop reads the clock again. If no deadline has passed it
waits again, so no timed event is handled early. For each timed event the server records its
deadline, when its wait returned and which pass handled it.

\subsection{Dispatch}\label{sec:dispatch}

\runin{In-process.} The classified connection stays on its worker, and its handler takes over.
TLS is terminated in the server.

\runin{Relay.} Relay dispatch exists on Linux only. The front opens a loopback connection to the
back-end port of the connection's class, writes any bytes replay has read, and copies both ways.
A side's end is passed on as a half-close, and the connection closes when both directions have
ended; a reset on one side closes the other by reset. For TLS the front never terminates. Its own
parser reassembles the ClientHello across records, up to \Bch{}, and routes by its server name
and ALPN list (pass-through). The relay copies through user-space buffers of \RelayBuf{} bytes
per direction, or with \texttt{splice}; rule~E chooses per backend.

\runin{Stub mode.} The back end of the comparisons with the proxies is the server in stub mode:
one port per class, and the least work that completes an exchange. HTTP/1.1 gets the handler's
response. On the TLS port the stub reads one TLS record, writes a body of \BodyBytes{} bytes and
closes; it never makes a handshake.

\subsection{Handlers and TLS}\label{sec:handlers}

Every handler is the same code in both modes, and each is minimal. HTTP/1.1 is the server's own
parser; it answers status \texttt{200} with a body of \BodyBytes{} bytes. A request line that
breaks the detector's grammar is closed without a response in both modes, so the modes agree on
invalid input. HTTP/2 runs through nghttp2 version \NghttpVersion~\cite{nghttp22026}. MQTT answers
CONNECT with CONNACK and PINGREQ with PINGRESP, and closes on DISCONNECT. SSH sends its identification line \SshLine{} on entry, reads the
client's line and closes. SMTP, reached only as a fallback or on its dedicated port, answers
\texttt{220}, then \texttt{250} to EHLO, \texttt{221} to QUIT and \texttt{500} to anything
else~\cite{klensin2008rfc5321}.

TLS runs through OpenSSL from its \OpensslSeries{} long-term support series, pinned at release
\OpensslVersion{} and built from source on both hosts~\cite{openssl2026}. The settings are the
same in the server, the back end and the generator. They are TLS~1.3 only, the cipher suite
\CipherSuite, the signature scheme \SigScheme{} and the group \GroupX{} only. There are no session
tickets (\NoTickets), no session cache and no early data. The cost cells offer ALPN
\texttt{http/1.1}. The loop owns the socket on every backend and gives OpenSSL its bytes through a
BIO of the worker's own. io\_uring and IOCP complete operations that the library cannot issue
itself.

\subsection{Counters}\label{sec:counters}

Every build counts, per worker, in plain integers without atomics, in both modes. It counts
operations by kind: accept, receive, peek, zero-byte receive, send, \texttt{setsockopt} and the
check of rule (b). It also counts \texttt{io\_uring\_enter} calls, submissions by opcode and
\texttt{GetQueuedCompletionStatusEx} calls. Further counters hold the switches from peek to replay
and the payload bytes across the boundary between user and kernel space, by direction. Others hold
the bytes copied in user space, the wakeups per connection during detection, the outcomes per class
and the timed events above. The counters run in every arm, so they cost the same on
both sides of each contrast.

\subsection{Generators and the Holder}\label{sec:generators}

Three first-party programs drive and observe the server. \opgen{} generates load: scripted
exchanges per protocol, a new connection per exchange or kept-alive connections, closed and open
loop, and the time to the first response byte. It binds each source address with
\texttt{IP\_BIND\_ADDRESS\_NO\_PORT}, and each window takes the next unused block of
\KSrc{} addresses of \LoopbackNet{}, never reused within \TimeWaitS{}~s. Windows has no such
option, so on W \opgen{} sets \texttt{SO\_REUSE\_UNICASTPORT} before each bind, a reading the
revision log records. Two facts go with it, both from development runs on W. A port is named right
after the bind with or without the option, so the port is not shown to be chosen at connect. And
the generator's churn rate did not change without the option. \opcase{} runs each hard
case, and the openings of the footprint test, as a script of writes and gaps, one write per chunk
with \texttt{TCP\_NODELAY}. It records each write's time on the server's clock. \ophold{} accepts
connections and holds them without reading or polling them. It measures the kernel's memory for a
bare held socket.
